\centerline{\bf \Huge Вероятность и статистика}

\problem{\textbf{1}} По расписанию автобус должен быть на остановке в 9:32. Расположите следующие события в порядке возрастания их вероятностей.\\
$A=$ «Автобус прибыл на остановку не раньше чем в 9:28».\\
$B=$ «Время прибытия автобуса на остановку не более чем на 1 минуту отклоняется от расписания».\\
$C=$ «Время прибытия автобуса на остановку не более чем на 4 минуты отклоняется от расписания».\\
$D=$ «Автобус прибыл на остановку не раньше чем в 9:31 и не позже чем в 9:35».

\problem{\textbf{2}} При однократном бросании несимметричной монетки, решка выпадает с вероятностью 0,2. Расположите события по возрастанию вероятностей:\\
$A=$ при двух бросаниях решка выпадет ровно один раз;\\
$B=$ при трех бросаниях решка выпадет ровно один раз;\\
$C=$ при трех бросаниях решка выпадет ровно два раза;\\
$D=$ при трех бросаниях решка не выпадет ни разу.

\problem{\textbf{3}} Монету бросили 20 раз. Известно, что орел выпал 9 раз. Найдите вероятность того, что при десятом по счету броске выпала решка.

\problem{\textbf{4}} За круглый стол на 201 стул в случайном порядке рассаживаются 199 мальчиков и 2 девочки. Найдите вероятность того, что между девочками будет сидеть один мальчик.

\problem{\textbf{5}} Ковбой Джон попадает в муху на стене с вероятностью 0,9, если стреляет из пристрелянного револьвера. Если Джон стреляет из непристрелянного револьвера, то он попадает в муху с вероятностью 0,2. На столе лежит 10 револьверов, из них только 4 пристрелянные. Ковбой Джон видит на стене муху, наудачу хватает первый попавшийся револьвер и стреляет в муху. Найдите вероятность того, что Джон промахнётся.

\problem{\textbf{6}} В торговом центре два одинаковых автомата продают кофе. Обслуживание автоматов происходит по вечерам после закрытия центра. Известно, что вероятность события «К вечеру в первом автомате закончится кофе» равна 0,25. Такая же вероятность события «К вечеру во втором автомате закончится кофе». Вероятность того, что кофе к вечеру закончится в обоих автоматах, равна 0,15. Найдите вероятность того, что к вечеру кофе останется в обоих автоматах.\\

\centerline{\textbf{40 минут}}